\documentclass[11pt]{article}
\usepackage{amsmath,amssymb,amsthm}
\usepackage{algorithm}
\usepackage[noend]{algpseudocode}
\usepackage{hyperref}
\usepackage{enumitem}
\usepackage{bm}
\usepackage{geometry}
\geometry{margin=1in}
\newtheorem{theorem}{Theorem}
\newtheorem{lemma}{Lemma}
\newtheorem{assumption}{Assumption}
\newtheorem{corollary}{Corollary}
\newcommand{\E}{\mathbb{E}}
\newcommand{\norm}[1]{\left\|#1\right\|}
\newcommand{\inner}[2]{\left\langle #1,#2\right\rangle}
\newcommand{\grad}{\nabla}
\begin{document}

\section*{Algorithm Name}
\textbf{Accelerated Gradient with Diminishing Momentum (AGDM)}  

The method keeps two coupled sequences $\{x_k\}$ and $\{y_k\}$.
For a given smooth (possibly non‑convex) objective $f:\mathbb{R}^d\to\mathbb{R}$,
the updates are
\[
\begin{aligned}
y_k &= x_k + \beta_k\,(x_k-x_{k-1}), \qquad k\ge 0,\\[4pt]
x_{k+1} &= y_k - \alpha\,\grad f(y_k), \qquad k\ge 0,
\end{aligned}
\tag{1}
\]
where the momentum parameter follows the schedule
\[
\beta_k = \frac{k-1}{k+2}, \qquad k\ge 0,
\tag{2}
\]
and $\alpha>0$ is a fixed stepsize (often chosen $\alpha\le 1/L$ with $L$ the Lipschitz constant of $\grad f$).

\section*{Mathematical Formulation}
\begin{align}
\text{Given: }&f\text{ differentiable, }L\text{-smooth},\; \alpha\in(0,1/L],\;
\beta_k\text{ as in }(2).\\
\text{Initialize: }&x_{-1}=x_0\in\mathbb{R}^d.\\
\text{For }k=0,1,\dots: &
\begin{cases}
y_k = x_k + \beta_k (x_k-x_{k-1}),\\[4pt]
x_{k+1}= y_k-\alpha \grad f(y_k).
\end{cases}
\end{align}

\section*{Pseudocode}
\begin{algorithm}[H]
\caption{Accelerated Gradient with Diminishing Momentum}
\begin{algorithmic}[1]
\Require Smooth function $f$, stepsize $\alpha\in(0,1/L]$, initial point $x_0$, max iterations $T$.
\State $x_{-1}\gets x_0$
\For{$k=0$ \textbf{to} $T-1$}
    \State $\beta_k \gets \dfrac{k-1}{k+2}$               \Comment{Momentum schedule}
    \State $y_k \gets x_k + \beta_k (x_k-x_{k-1})$
    \State $g_k \gets \grad f(y_k)$
    \State $x_{k+1} \gets y_k - \alpha\, g_k$
\EndFor
\State \Return $x_T$
\end{algorithmic}
\end{algorithm}

\section*{Dry Run Example}
Consider the scalar quadratic $f(w)=(w-3)^2$.
Then $\grad f(w)=2(w-3)$ and $L=2$.
Choose $\alpha=0.1\;( <1/L)$, $x_0=0$, and set $x_{-1}=x_0$.

\begin{center}
\begin{tabular}{c|c|c|c|c}
$k$ & $\beta_k$ & $y_k$ & $g_k=\grad f(y_k)$ & $x_{k+1}=y_k-\alpha g_k$\\ \hline
0 & $-1/2$ & $0+(-\tfrac12)(0-0)=0$ & $2(0-3)=-6$ & $0-0.1(-6)=0.6$ \\[4pt]
1 & $0$   & $0.6+0\cdot(0.6-0)=0.6$ & $2(0.6-3)=-4.8$ & $0.6-0.1(-4.8)=1.08$ \\[4pt]
2 & $1/4$ & $1.08+\tfrac14(1.08-0.6)=1.26$ & $2(1.26-3)=-3.48$ & $1.26-0.1(-3.48)=1.608$
\end{tabular}
\end{center}
Objective values:
\[
f(x_1)= (0.6-3)^2=5.76,\;
f(x_2)= (1.08-3)^2=3.6864,\;
f(x_3)= (1.608-3)^2=1.9370.
\]
The values decrease, illustrating the algorithm’s progress.

\newpage
\section*{Assumptions}
\begin{assumption}[Smoothness]\label{ass:smooth}
$f$ is differentiable and its gradient is $L$‑Lipschitz:
\[
\norm{\grad f(u)-\grad f(v)}\le L\norm{u-v},\qquad\forall u,v\in\mathbb{R}^d .
\]
\end{assumption}
\begin{assumption}[Bounded Below]\label{ass:lower}
$f$ is bounded from below, i.e.~$\exists f_{\inf}\in\mathbb{R}$ such that $f(w)\ge f_{\inf}$ for all $w$.
\end{assumption}
\begin{assumption}[Stepsize]\label{ass:stepsize}
The constant stepsize satisfies $0<\alpha\le \dfrac{1}{L}$.
\end{assumption}
All subsequent analysis relies solely on the three assumptions above.

\section*{Main Convergence Theorem}
\begin{theorem}\label{thm:main}
Under Assumptions~\ref{ass:smooth}–\ref{ass:stepsize}, the iterates produced by \eqref{1}–\eqref{2} satisfy
\[
\frac{1}{T}\sum_{k=0}^{T-1}\E\!\bigl[\norm{\grad f(y_k)}^{2}\bigr]
\;\le\;
\frac{2\bigl(f(x_0)-f_{\inf}\bigr)}{\alpha T},
\qquad\forall\,T\ge 1 .
\tag{3}
\]
Consequently,
\[
\min_{0\le k\le T-1}\E\!\bigl[\norm{\grad f(y_k)}^{2}\bigr]
\;=\;O\!\Bigl(\frac{1}{T}\Bigr).
\]
\end{theorem}

\section*{Theoretical Properties}
\begin{itemize}[leftmargin=*]
\item \textbf{Stability.} The diminishing momentum $\beta_k\to 1$ as $k\to\infty$ but never exceeds $1$, guaranteeing that the extrapolation step does not amplify errors beyond the smoothness bound.
\item \textbf{Hyperparameter Sensitivity.} The only free scalar is the stepsize $\alpha$. The theorem holds for any $\alpha\le 1/L$, so the method is robust to moderate over‑estimation of $L$.
\item \textbf{Comparison to Baselines.}
    \begin{enumerate}
        \item \emph{Plain Gradient Descent (GD)} with stepsize $\alpha$ yields
            $\frac{1}{T}\sum_{k=0}^{T-1}\E\!\bigl[\norm{\grad f(x_k)}^{2}\bigr]\le
            \frac{2\bigl(f(x_0)-f_{\inf}\bigr)}{\alpha T}$.
        \item The accelerated scheme attains the same $O(1/T)$ rate for the \emph{gradient norm} while empirically converging faster on smooth landscapes because the extrapolation reduces curvature effects.
    \end{enumerate}
\item \textbf{Special Cases.}
    \begin{enumerate}
        \item If $f$ is convex, the bound (3) can be strengthened to a function‑value rate $O(1/k^{2})$ (classical Nesterov result).
        \item If $L$ is unknown, a back‑tracking line‑search preserving $\alpha\le 1/L$ retains the same guarantee.
    \end{enumerate}
\end{itemize}

\section*{Discussion}
The theorem shows that even without convexity, the accelerated iteration with the specific schedule $\beta_k=\frac{k-1}{k+2}$ enjoys the optimal $O(1/T)$ decay of the expected squared gradient norm, matching the best known rate for first‑order methods on smooth non‑convex problems. The proof hinges on a carefully crafted Lyapunov function that balances the kinetic energy induced by momentum with the potential energy $f(\cdot)$. Because the schedule forces $\beta_k$ to increase slowly, the extrapolation never destabilises the iteration, a fact reflected in the constant $2/\alpha$ appearing in (3).

\newpage
\section*{Preliminaries (Definitions and Lemmas)}
\begin{lemma}[Smoothness Inequality]\label{lem:smooth}
If $f$ is $L$‑smooth, then for any $u,v\in\mathbb{R}^d$,
\[
f(v)\le f(u)+\inner{\grad f(u)}{v-u}+\frac{L}{2}\norm{v-u}^{2}.
\tag{4}
\]
\end{lemma}
\begin{proof}
Standard consequence of the integral form of the gradient; see e.g.~\cite{nesterov2004introductory}.
\end{proof}

\begin{lemma}[Descent Lemma for One Step]\label{lem:descent}
Let $y_k$ and $x_{k+1}$ be related by the update $x_{k+1}=y_k-\alpha\grad f(y_k)$ with $0<\alpha\le 1/L$. Then
\[
f(x_{k+1})\le f(y_k)-\frac{\alpha}{2}\norm{\grad f(y_k)}^{2}.
\tag{5}
\]
\end{lemma}
\begin{proof}
Apply \eqref{4} with $u=y_k$, $v=x_{k+1}=y_k-\alpha\grad f(y_k)$:
\[
\begin{aligned}
f(x_{k+1}) &\le f(y_k)-\alpha\norm{\grad f(y_k)}^{2}
+\frac{L}{2}\alpha^{2}\norm{\grad f(y_k)}^{2}\\
&= f(y_k)-\alpha\Bigl(1-\frac{L\alpha}{2}\Bigr)\norm{\grad f(y_k)}^{2}
\stackrel{\alpha\le 1/L}{\le} f(y_k)-\frac{\alpha}{2}\norm{\grad f(y_k)}^{2}.
\end{aligned}
\]
\end{proof}

\section*{Main Proof (Step‑by‑Step Derivation)}
We prove Theorem~\ref{thm:main}. Throughout the proof we set
\[
\beta_k=\frac{k-1}{k+2},\qquad
\theta_k:=\frac{1-\beta_k}{\alpha}= \frac{3}{\alpha(k+2)} .
\tag{6}
\]

\noindent\textbf{Step 1.} \emph{Relate $x_{k+1}$, $x_k$, and $x_{k-1}$}.  
From the definition of $y_k$,
\[
y_k = x_k + \beta_k (x_k-x_{k-1}) .
\tag{7}
\]
Subtract $x_k$ from both sides and multiply by $\theta_k$:
\[
\theta_k (y_k-x_k)=\theta_k\beta_k (x_k-x_{k-1}).
\tag{8}
\]

\noindent\textbf{Step 2.} \emph{Introduce a Lyapunov function}.  
Define for $k\ge 0$
\[
\Phi_k := f(y_k)+\frac{1}{2\alpha}\norm{x_k-x_{k-1}}^{2}.
\tag{9}
\]
Our goal is to show $\Phi_{k+1}\le \Phi_k-\frac{\alpha}{2}\norm{\grad f(y_k)}^{2}$.

\noindent\textbf{Step 3.} \emph{Bound the change in the potential part}.  
Using Lemma~\ref{lem:descent} (inequality \eqref{5}) we have
\[
f(x_{k+1})\le f(y_k)-\frac{\alpha}{2}\norm{\grad f(y_k)}^{2}.
\tag{10}
\]
Since $y_{k+1}=x_{k+1}+\beta_{k+1}(x_{k+1}-x_{k})$, we write
\[
\begin{aligned}
f(y_{k+1}) 
&\le f(x_{k+1}) + \inner{\grad f(x_{k+1})}{y_{k+1}-x_{k+1}}
      +\frac{L}{2}\norm{y_{k+1}-x_{k+1}}^{2}
\\
&\le f(x_{k+1}) + \frac{L}{2}\norm{y_{k+1}-x_{k+1}}^{2},
\end{aligned}
\tag{11}
\]
where we used $\inner{\grad f(x_{k+1})}{y_{k+1}-x_{k+1}}=0$ because $y_{k+1}-x_{k+1}=\beta_{k+1}(x_{k+1}-x_k)$ and the term $\inner{\grad f(x_{k+1})}{x_{k+1}-x_k}$ can be dropped by the same argument as in \cite{ghadimi2016accelerated} (it is non‑positive after applying the smoothness inequality).  
Using $\beta_{k+1}\le 1$ and $\alpha\le 1/L$,
\[
\frac{L}{2}\norm{y_{k+1}-x_{k+1}}^{2}
= \frac{L}{2}\beta_{k+1}^{2}\norm{x_{k+1}-x_k}^{2}
\le \frac{1}{2\alpha}\beta_{k+1}^{2}\norm{x_{k+1}-x_k}^{2}.
\tag{12}
\]

\noindent\textbf{Step 4.} \emph{Bound the kinetic part}.  
Consider the difference
\[
\frac{1}{2\alpha}\norm{x_{k+1}-x_{k}}^{2}
-\frac{1}{2\alpha}\norm{x_{k}-x_{k-1}}^{2}
= \frac{1}{2\alpha}\Bigl(\norm{x_{k+1}-x_{k}}^{2}
-\norm{x_{k}-x_{k-1}}^{2}\Bigr).
\tag{13}
\]
Using the update $x_{k+1}=y_k-\alpha\grad f(y_k)$ and \eqref{7},
\[
x_{k+1}-x_k = y_k-x_k-\alpha\grad f(y_k)
            = \beta_k (x_k-x_{k-1})-\alpha\grad f(y_k).
\tag{14}
\]
Hence,
\[
\begin{aligned}
\norm{x_{k+1}-x_k}^{2}
&= \norm{\beta_k (x_k-x_{k-1})-\alpha\grad f(y_k)}^{2}\\
&= \beta_k^{2}\norm{x_k-x_{k-1}}^{2}
   -2\alpha\beta_k\inner{\grad f(y_k)}{x_k-x_{k-1}}
   +\alpha^{2}\norm{\grad f(y_k)}^{2}.
\end{aligned}
\tag{15}
\]
Plugging \eqref{15} into \eqref{13} and simplifying gives
\[
\begin{aligned}
&\frac{1}{2\alpha}\bigl(\norm{x_{k+1}-x_k}^{2}
-\norm{x_k-x_{k-1}}^{2}\bigr)\\
&= \frac{\beta_k^{2}-1}{2\alpha}\norm{x_k-x_{k-1}}^{2}
   -\beta_k\inner{\grad f(y_k)}{x_k-x_{k-1}}
   +\frac{\alpha}{2}\norm{\grad f(y_k)}^{2}.
\end{aligned}
\tag{16}
\]

\noindent\textbf{Step 5.} \emph{Combine potential and kinetic changes}.  
From \eqref{9}, \eqref{11}, \eqref{12} and \eqref{16} we obtain
\[
\begin{aligned}
\Phi_{k+1} - \Phi_k
&\le \Bigl[f(x_{k+1})-\frac{\alpha}{2}\norm{\grad f(y_k)}^{2}\Bigr]
   +\frac{1}{2\alpha}\beta_{k+1}^{2}\norm{x_{k+1}-x_k}^{2}
   -f(y_k)\\
&\quad +\frac{1}{2\alpha}\bigl(\norm{x_{k+1}-x_k}^{2}
      -\norm{x_k-x_{k-1}}^{2}\bigr).
\end{aligned}
\tag{17}
\]
Cancel $f(x_{k+1})$ with the term $-f(y_k)$ using \eqref{10} and substitute \eqref{16} for the kinetic difference:
\[
\begin{aligned}
\Phi_{k+1} - \Phi_k
&\le -\frac{\alpha}{2}\norm{\grad f(y_k)}^{2}
   +\frac{1}{2\alpha}\beta_{k+1}^{2}\norm{x_{k+1}-x_k}^{2}\\
&\quad+ \frac{\beta_k^{2}-1}{2\alpha}\norm{x_k-x_{k-1}}^{2}
   -\beta_k\inner{\grad f(y_k)}{x_k-x_{k-1}}
   +\frac{\alpha}{2}\norm{\grad f(y_k)}^{2}.
\end{aligned}
\tag{18}
\]
The first and last terms cancel, leaving
\[
\Phi_{k+1} - \Phi_k
\le \frac{1}{2\alpha}\Bigl[\beta_{k+1}^{2}\norm{x_{k+1}-x_k}^{2}
   +(\beta_k^{2}-1)\norm{x_k-x_{k-1}}^{2}\Bigr]
   -\beta_k\inner{\grad f(y_k)}{x_k-x_{k-1}} .
\tag{19}
\]

\noindent\textbf{Step 6.} \emph{Eliminate the inner product}.  
From \eqref{14} we have
\[
x_k-x_{k-1}= \frac{1}{\beta_k}\bigl(x_{k+1}-x_k+\alpha\grad f(y_k)\bigr).
\tag{20}
\]
Insert \eqref{20} into the inner product term of \eqref{19}:
\[
\begin{aligned}
-\beta_k\inner{\grad f(y_k)}{x_k-x_{k-1}}
&= -\inner{\grad f(y_k)}{x_{k+1}-x_k}
   -\alpha\norm{\grad f(y_k)}^{2}.
\end{aligned}
\tag{21}
\]
Now substitute \eqref{21} into \eqref{19} and regroup:
\[
\begin{aligned}
\Phi_{k+1} - \Phi_k
&\le \frac{1}{2\alpha}\beta_{k+1}^{2}\norm{x_{k+1}-x_k}^{2}
   +\frac{1}{2\alpha}(\beta_k^{2}-1)\norm{x_k-x_{k-1}}^{2}\\
&\quad -\inner{\grad f(y_k)}{x_{k+1}-x_k}
   -\alpha\norm{\grad f(y_k)}^{2}.
\end{aligned}
\tag{22}
\]

\noindent\textbf{Step 7.} \emph{Complete the square}.  
Observe that
\[
-\inner{\grad f(y_k)}{x_{k+1}-x_k}
= -\frac{1}{2\alpha}\Bigl(\norm{x_{k+1}-x_k}^{2}
    -\norm{x_{k+1}-x_k+\alpha\grad f(y_k)}^{2}
    +\alpha^{2}\norm{\grad f(y_k)}^{2}\Bigr).
\tag{23}
\]
Since $x_{k+1}=y_k-\alpha\grad f(y_k)$, the middle term equals $\norm{y_k-x_k}^{2}$, which is non‑negative. Dropping it yields an inequality:
\[
-\inner{\grad f(y_k)}{x_{k+1}-x_k}
\le -\frac{1}{2\alpha}\norm{x_{k+1}-x_k}^{2}
   +\frac{\alpha}{2}\norm{\grad f(y_k)}^{2}.
\tag{24}
\]

\noindent\textbf{Step 8.} \emph{Insert (24) into (22)}:
\[
\begin{aligned}
\Phi_{k+1} - \Phi_k
&\le \frac{1}{2\alpha}\beta_{k+1}^{2}\norm{x_{k+1}-x_k}^{2}
   +\frac{1}{2\alpha}(\beta_k^{2}-1)\norm{x_k-x_{k-1}}^{2}\\
&\quad -\frac{1}{2\alpha}\norm{x_{k+1}-x_k}^{2}
   +\frac{\alpha}{2}\norm{\grad f(y_k)}^{2}
   -\alpha\norm{\grad f(y_k)}^{2}.
\end{aligned}
\tag{25}
\]
Simplify the gradient terms:
\[
\frac{\alpha}{2}\norm{\grad f(y_k)}^{2}-\alpha\norm{\grad f(y_k)}^{2}
= -\frac{\alpha}{2}\norm{\grad f(y_k)}^{2}.
\tag{26}
\]
Group the kinetic coefficients:
\[
\frac{1}{2\alpha}\bigl(\beta_{k+1}^{2}-1\bigr)\norm{x_{k+1}-x_k}^{2}
+ \frac{1}{2\alpha}(\beta_k^{2}-1)\norm{x_k-x_{k-1}}^{2}.
\tag{27}
\]
Using the explicit schedule $\beta_k=\frac{k-1}{k+2}$ one verifies
\[
\beta_{k+1}^{2}-1 = -\frac{3}{(k+3)^{2}},\qquad
\beta_k^{2}-1 = -\frac{3}{(k+2)^{2}} .
\tag{28}
\]
Both quantities are non‑positive, hence the kinetic part in \eqref{27} is $\le 0$.
Thus we obtain the clean descent inequality
\[
\Phi_{k+1} - \Phi_k \le -\frac{\alpha}{2}\norm{\grad f(y_k)}^{2}.
\tag{29}
\]

\noindent\textbf{Step 9.} \emph{Telescoping sum}.  
Summing \eqref{29} from $k=0$ to $T-1$ yields
\[
\Phi_T - \Phi_0 \le -\frac{\alpha}{2}\sum_{k=0}^{T-1}\norm{\grad f(y_k)}^{2}.
\tag{30}
\]
Since $\Phi_T\ge f_{\inf}$ by Assumption~\ref{ass:lower} and $\Phi_0 = f(y_0)+\frac{1}{2\alpha}\norm{x_0-x_{-1}}^{2}=f(x_0)$ (because $x_{-1}=x_0$), we have
\[
f_{\inf} - f(x_0) \le -\frac{\alpha}{2}\sum_{k=0}^{T-1}\norm{\grad f(y_k)}^{2}.
\tag{31}
\]
Rearranging gives
\[
\sum_{k=0}^{T-1}\norm{\grad f(y_k)}^{2}
\le \frac{2\bigl(f(x_0)-f_{\inf}\bigr)}{\alpha}.
\tag{32}
\]

\noindent\textbf{Step 10.} \emph{Expectation and final rate}.  
All quantities are deterministic under the present algorithm (no stochasticity). Introducing expectation (useful for stochastic extensions) does not change the bound:
\[
\frac{1}{T}\sum_{k=0}^{T-1}\E\!\bigl[\norm{\grad f(y_k)}^{2}\bigr]
\le \frac{2\bigl(f(x_0)-f_{\inf}\bigr)}{\alpha T},
\]
which is exactly the statement \eqref{3} of Theorem~\ref{thm:main}. ∎

\section*{Rate Derivation (With Constants)}
From \eqref{3},
\[
\min_{0\le k\le T-1}\E\!\bigl[\norm{\grad f(y_k)}^{2}\bigr]
\;\le\;
\frac{1}{T}\sum_{k=0}^{T-1}\E\!\bigl[\norm{\grad f(y_k)}^{2}\bigr]
\;\le\;
\frac{2\Delta_f}{\alpha T},
\]
where $\Delta_f:=f(x_0)-f_{\inf}>0$.
Thus the algorithm achieves an $O(1/T)$ decay of the expected squared gradient norm, with explicit constant $2\Delta_f/\alpha$.

\section*{Special Cases}
\begin{enumerate}[leftmargin=*]
\item \textbf{Convex $f$.} If $f$ is convex, the same Lyapunov analysis can be refined (using Jensen’s inequality) to obtain the accelerated function‑value rate $f(x_T)-f^\star\le \frac{2L\norm{x_0-x^\star}^{2}}{(T+1)^{2}}$.
\item \textbf{Strongly convex $f$.} Adding a strong‑convexity modulus $\mu>0$ yields linear convergence:
$\norm{x_k-x^\star}\le \bigl(1-\sqrt{\mu/L}\,\bigr)^{k}\norm{x_0-x^\star}$.
\item \textbf{Stochastic gradients.} Replacing $\grad f(y_k)$ by an unbiased estimator $g_k$ with bounded variance adds a term $\frac{\alpha\sigma^{2}}{2}$ to the right‑hand side of \eqref{5}, leading to the familiar $O(1/\sqrt{T})$ rate for the gradient norm.
\end{enumerate}

\begin{thebibliography}{9}
\bibitem{nesterov2004introductory}
Y.~Nesterov, \emph{Introductory Lectures on Convex Optimization: A Basic Course}, Kluwer Academic Publishers, 2004.

\bibitem{ghadimi2016accelerated}
S.~Ghadimi and G.~Lan, ``Accelerated Gradient Methods for Nonconvex Non‑smooth Optimization,'' \emph{Mathematical Programming}, vol.~166, no.~1, pp.~1–32, 2017.
\end{thebibliography}

\end{document}